\documentclass[11pt,a4paper]{article}
\usepackage[T1]{fontenc}
\usepackage{lmodern}
\usepackage[margin=25mm,headheight=14pt]{geometry}
\usepackage{amsmath,amssymb,mathtools,bm}
\usepackage{booktabs,array,enumitem,microtype}
\usepackage[dvipsnames]{xcolor}
\usepackage{hyperref,fancyhdr}
\hypersetup{colorlinks=true,linkcolor=MidnightBlue,urlcolor=MidnightBlue,
  pdftitle={Learning conditional loading priors in cEBMF},
  pdfauthor={Methods note for the cEBMF project}}
\pagestyle{fancy}\fancyhf{}
\fancyhead[L]{\small Conditional loading priors in cEBMF}
\fancyhead[R]{\small Derivation}
\fancyfoot[C]{\thepage}
\setlength{\parindent}{0pt}\setlength{\parskip}{5pt}
\setlist{itemsep=3pt,topsep=4pt,leftmargin=*}
\allowdisplaybreaks[1]
\newcommand{\E}{\mathbb E}
\newcommand{\N}{\mathcal N}
\newcommand{\LL}{\mathcal L}
\newcommand{\pa}{\operatorname{pa}}
\newcommand{\ch}{\operatorname{ch}}
\newcommand{\Cov}{\operatorname{Cov}}
\newcommand{\sigmoid}{\operatorname{sigmoid}}
\newcommand{\logit}{\operatorname{logit}}
\newcommand{\Normal}[3]{\varphi(#1;#2,#3)}
\newcommand{\note}[1]{\par\smallskip\noindent\colorbox{blue!5}{\parbox{\dimexpr\linewidth-2\fboxsep\relax}{#1}}\par\smallskip}
\title{\textbf{Learning conditional loading priors in cEBMF}\\[7pt]
  \Large $p_\theta(L_k\mid L_{<k})$ and a learning algorithm\\
  that reduces to ordinary cEBMF}
\author{Methods note for the cEBMF project}
\date{16 September 2026}

\begin{document}
\maketitle
\thispagestyle{fancy}

\note{\textbf{The modeling objective.} Learn the conditional distributions
\(p_\theta(L_k\mid L_{<k})\), and hence the dependent joint prior
\(p_\theta(L)=\prod_kp_\theta(L_k\mid L_{<k})\), while estimating the latent
loadings from the observed matrix. Keep an ordinary learned ash prior on
\(F\). The Gaussian-mixture update is implemented with numerical integration.
This note distinguishes its exact variational derivation from its computational approximations.}

\section{The model we want to learn}
Let \(Y\in\mathbb R^{n\times p}\), \(L\in\mathbb R^{n\times K}\), and
\(F\in\mathbb R^{p\times K}\). Here \(L_k=L_{:k}\) means the \emph{kth
loading column}, and \(L_{<k}=L_{:,1:(k-1)}\). The observation model is
\begin{equation}
Y_{ij}\mid L,F,\tau\sim
\N\!\left(\sum_{k=1}^KL_{ik}F_{jk},\tau_{ij}^{-1}\right),
\qquad (i,j)\in\Omega.
\label{eq:likelihood}
\end{equation}
The desired autoregressive loading prior is
\begin{equation}
\boxed{p_\theta(L\mid X)=\prod_{k=1}^Kp_{\theta_k}(L_k\mid L_{<k},X).}
\label{eq:desired}
\end{equation}
As in the existing conditional-prior networks, we use the rowwise form
\begin{equation}
p_{\theta_k}(L_k\mid L_{<k},X)
=\prod_{i=1}^ng_{\theta_k}(L_{ik}\mid X_i,L_{i,<k}).
\label{eq:rows}
\end{equation}
The parameters \(\theta_k\) are shared across rows. If fixed covariates are
unavailable, omit \(X_i\). Conditional independence across rows in
\eqref{eq:rows} is a specified architecture; it is not independence across
loading columns.

For example, a CGB network learns
\[
g_{\theta_k}(dv\mid P)=\pi_{k0}(P)\delta_0(dv)
+[1-\pi_{k0}(P)]\N(dv;\mu_k,s_k^2),\quad P=(X_i,L_{i,<k}).
\]
Earlier loadings predict whether a later loading is active. An EMDN can
also let earlier loadings predict component means and variances. Multiplying
these normalized, acyclic conditionals defines a normalized joint prior.
The ordering permits dependencies; it does not uniquely identify a tree.

\clearpage
\section{Prior dependence and posterior approximation}
The earlier loadings are latent, not observed covariates. A parent loading
therefore receives information both from the matrix likelihood and from
later loading-prior terms. For two loading coordinates in one row,
\[
p(l_1,l_2)=g_1(l_1)g_2(l_2\mid l_1)
\quad\Longrightarrow\quad
p(l_1\mid l_2,Y,F)\propto p(Y\mid l_1,l_2,F)g_1(l_1)g_2(l_2\mid l_1).
\]
This posterior feedback follows from the requested joint model itself.
It does not introduce a reverse-direction generative prior.

\subsection{Mixture labels and the variational family}
For clarity, represent a loading mixture by a label \(Z_{ik}\):
\begin{align}
p(Z_{ik}=0,L_{ik}\in dv\mid P_{ik})
  &=\pi_{k0}(P_{ik})\delta_0(dv),\notag\\
p(Z_{ik}=h,L_{ik}\in dv\mid P_{ik})
  &=\pi_{kh}(P_{ik})\Normal{v}{\mu_{kh}(P_{ik})}{s_{kh}^2(P_{ik})}\,dv,
  \quad h>0.
\label{eq:mixture}
\end{align}
Here \(P_{ik}\) contains the fixed covariates and the selected earlier
parents. These labels only identify existing mixture components. They
are not an additional representation or upper latent layer.

Use the variational approximation
\begin{equation}
q(L,Z,F)=\prod_{i,k}q_{ik}(L_{ik},Z_{ik})\prod_{j,k}q^F_{jk}(F_{jk}).
\label{eq:q}
\end{equation}
Each value stays dependent on its own label. The generative prior
\eqref{eq:desired} remains dependent even though this posterior
approximation factorizes across loading coordinates. Thus this method
\emph{learns the requested conditional prior}, but does not represent
posterior covariance between distinct loading coordinates. Section~\ref{sec:structured}
states the alternative if that covariance is also required.

For overlapping Gaussian components, augmenting the labels before making
the mean-field approximation can give a different approximation from
marginalizing them first. Both describe the same generative marginal model.
With fixed parents and no children, the local optimum is the usual
normal-means posterior, so augmentation does not spoil the cEBMF limit.

\subsection{A single objective}
Let \(g^F_{\eta_k}\) be the ordinary ash prior on feature column \(k\), with
mixture weights estimated during fitting. The unpenalized ELBO is
\begin{align}
\LL(q,\theta,\eta,\tau)
={}&\E_q\log p(Y\mid L,F,\tau)+H(q)\notag\\
 &+\sum_{i,k}\E_q\log g_{\theta_k}(L_{ik},Z_{ik}\mid P_{ik})
 +\sum_{j,k}\E_q\log g^F_{\eta_k}(F_{jk}).
\label{eq:elbo}
\end{align}
For the loading terms, densities and entropy are taken with respect to a
label-specific reference measure: unit mass at zero for the spike and
Lebesgue measure for continuous components. Gaussian densities are not
added to spike probability mass at zero.

\clearpage
\section{Deriving a loading-coordinate update}
\subsection{The original cEBMF likelihood calculation}
Write \(w_{ij}=M_{ij}\tau_{ij}\), where \(M\) is the observation mask. Let bars
denote first posterior moments, and define
\begin{align}
\bar R_{ij,-k}&=Y_{ij}-\sum_{h\ne k}\bar L_{ih}\bar F_{jh},\notag\\
a_{ik}&=\sum_jw_{ij}\E_q[F_{jk}^2],&
b_{ik}&=\sum_jw_{ij}\bar R_{ij,-k}\bar F_{jk}.
\label{eq:stats}
\end{align}
All sums involving observed Y are over observed entries. Expanding the
expected squared error while holding the other variational factors fixed
gives the loading-dependent likelihood term
\begin{equation}
\E_{q_{-ik}}\log p(Y\mid L,F,\tau)
=-\tfrac12a_{ik}v^2+b_{ik}v+\text{constant},\qquad v=L_{ik}.
\label{eq:quadratic}
\end{equation}
The denominator uses \(\E[F_{jk}^2]\), not \((\E F_{jk})^2\), as in
ordinary cEBMF. F need not be sampled.

\subsection{The own-prior and child terms}
Define
\begin{equation}
A_{ik}(v,z;\theta_k)=
\E_{q(P_{ik})}\log g_{\theta_k}(v,z\mid P_{ik}).
\label{eq:own}
\end{equation}
Only latent parents are averaged; the external covariates are fixed.
Changing v also changes every direct child's prior. Their contribution is
\begin{equation}
D_{ik}(v)=\sum_{c\in\ch(k)}
\E_{q(L_{ic},Z_{ic})q(P_{ic}\setminus L_{ik})}
\log g_{\theta_c}(L_{ic},Z_{ic}\mid P_{ic}[v]).
\label{eq:child}
\end{equation}
Under the total ordering, every later coordinate is a possible direct
child; structurally omitted network inputs create no edge. Holding the
other q factors and parameters fixed, the variational coordinate is
\begin{equation}
\boxed{q^*_{ik}(v,z)\propto
\exp\{-a_{ik}v^2/2+b_{ik}v+A_{ik}(v,z;\theta_k)+D_{ik}(v)\}.}
\label{eq:qstar}
\end{equation}
This follows by maximizing the terms of \eqref{eq:elbo} involving \(q_{ik}\):
expected log joint plus its entropy.

\note{\textbf{Two necessary corrections to plug-in conditioning.}
Average the \emph{log prior} over uncertain parents and include the
children's prior terms. In general,
\(\E\log g(v\mid P)\ne\log g(v\mid\E P)\) and
\(\E\log g(v\mid P)\ne\log\E g(v\mid P)\).
Using fitted earlier loadings as fixed inputs does not perform these steps.}

\clearpage
\section{Learning the conditional prior and posterior together}
Write the exponent in \eqref{eq:qstar} as
\(\Psi_{ik}(v,z;\theta_k)\), and define its normalizer
\begin{equation}
\mathcal Z_{ik}(\theta_k)=\sum_z\int
\exp\{\Psi_{ik}(v,z;\theta_k)\}\,d\nu_z(v).
\label{eq:normalizer}
\end{equation}
The network parameters are distinct for different loading columns.
Consequently \(D\) is fixed when optimizing \(\theta_k\), even though it varies
with the candidate loading \(v\). The Gibbs variational identity gives
\begin{equation}
\max_{q_{ik}}\{\E_{q_{ik}}\Psi_{ik}+H(q_{ik})\}
=\log\mathcal Z_{ik}(\theta_k).
\label{eq:profile}
\end{equation}
The joint prior-and-posterior update is therefore
\begin{equation}
\boxed{\theta_k^{\rm new}\in\arg\max_{\theta_k}\sum_i\log\mathcal Z_{ik}(\theta_k),
\qquad q_{ik}^{\rm new}=q^*_{ik}(\cdot;\theta_k^{\rm new}).}
\label{eq:learning}
\end{equation}
The neural prior is learned from all rows using the same ELBO as the
loading inference. An exact block optimizer cannot decrease that ELBO.
Approximate integration and finite neural optimization require their own
accuracy controls; they do not inherit an automatic monotonicity guarantee.

\subsection{Proof of the reduction to ordinary cEBMF}
With no latent-covariate edges, there are neither uncertain parents nor
child messages:
\[
A_{ik}(v,z;\theta_k)=\log g_{\theta_k}(v,z\mid X_i),\qquad D_{ik}(v)=0.
\]
For \(a=a_{ik}>0\), set \(x=b_{ik}/a\) and \(s^2=1/a\). Then
\[
e^{-av^2/2+bv}=\sqrt{2\pi/a}\,e^{b^2/(2a)}\Normal{x}{v}{1/a}.
\]
After summing over the component label,
\begin{align}
\log\mathcal Z_{ik}(\theta_k)
={}&\log\left\{\int\Normal{x}{v}{s^2}\,
g_{\theta_k}(dv\mid X_i)\right\}\notag\\
 &+\tfrac12\log(2\pi/a)+b^2/(2a).
\label{eq:reduction}
\end{align}
The last two terms do not depend on \(\theta_k\). Thus \eqref{eq:learning} is
\textbf{exactly the original cEBNM marginal-likelihood fit}, and
\eqref{eq:qstar} gives its posterior. Feature and noise updates below also
reduce to their ordinary cEBMF forms.

In code, use this analytic limit and call the existing cEBNM solver with
the same options, initialization and update schedule. A single
complete-data EM M-step, even using infinitely many exact posterior draws,
is not generally the same update as this marginal-likelihood fit.

For \(a=b=0\) use \eqref{eq:qstar} directly. Also, do not renormalize \(\exp(A)\)
as an ordinary prior and discard its normalizer: with uncertain parents
that normalizer generally depends on \(\theta_k\).

\clearpage
\section{Gaussian mixtures and exact zero spikes}
Suppress the row and loading indices. For each continuous component h,
the expected log own prior is quadratic in v. Define
\begin{align}
\alpha_h&=a+\E_P[s_h(P)^{-2}],\notag\\
\beta_h&=b+\E_P[\mu_h(P)s_h(P)^{-2}],\notag\\
\gamma_h&=\E_P\log\pi_h(P)-\tfrac12\E_P\log(2\pi s_h(P)^2)
-\tfrac12\E_P[\mu_h(P)^2s_h(P)^{-2}].
\label{eq:abc}
\end{align}
Assume these expectations are finite and \(\alpha_h>0\). Put
\(m_h=\beta_h/\alpha_h\), \(V_h=1/\alpha_h\), and
\begin{equation}
I_{rh}=\int v^r\Normal{v}{m_h}{V_h}e^{D(v)}dv,
\qquad r=0,1,2.
\label{eq:integrals}
\end{equation}
Completing the square gives unnormalized component masses
\begin{align}
W_h&=\exp\{\gamma_h+\beta_h^2/(2\alpha_h)\}
\sqrt{2\pi/\alpha_h}\,I_{0h},\quad h>0,\notag\\
W_0&=\exp\{\E_P\log\pi_0(P)+D(0)\}.
\label{eq:weights}
\end{align}
The zero-spike likelihood factor is one in the natural-parameter form
\(e^{-av^2/2+bv}\). There is no omitted Gaussian evidence constant in \(W_0\).
Normalize with \(r_h=W_h/\sum_tW_t\). Then
\begin{align}
q(v\mid z=h)&=\frac{\Normal{v}{m_h}{V_h}e^{D(v)}}{I_{0h}},\quad h>0,
\notag\\
\E_qv&=\sum_{h>0}r_hI_{1h}/I_{0h},&
\E_qv^2&=\sum_{h>0}r_hI_{2h}/I_{0h}.
\label{eq:moments}
\end{align}
Omit the spike for an unspiked EMDN. Compute component masses in log
space. With deterministic parents and \(D=0\) these are the ordinary analytic
Gaussian-mixture normal-means formulas.

\subsection{How to calculate the child message}
For child c, let \(r_{ch}=q(Z_c=h)\),
\(u_{ch}=\E(L_c\mid Z_c=h)\), and
\(T_{ch}=\E(L_c^2\mid Z_c=h)\). A Gaussian component contributes
\begin{equation}
J_{ch}(P)=-\tfrac12\log(2\pi s_{ch}^2(P))
-\frac{T_{ch}-2\mu_{ch}(P)u_{ch}+\mu_{ch}(P)^2}{2s_{ch}^2(P)}.
\label{eq:childdensity}
\end{equation}
Consequently,
\begin{equation}
D(v)=\sum_{c\in\ch(k)}\sum_h r_{ch}
\E_{P_c\setminus v}\left[\log\pi_{ch}(P_c[v])+J_{ch}(P_c[v])\right],
\label{eq:Dmoments}
\end{equation}
with \(J_{c0}=0\) for the fixed zero atom. For CGB with a global slab, J is
constant in v and drops from the v-dependent tilt. For EMDN, its means
and variances depend on parents, so J must remain. Child labels and values
can be integrated using their probabilities and moments rather than sampled.

\clearpage
\section{A two-loading example of the desired dependency}
As a transparent illustration, suppose the loadings are strictly binary:
\begin{equation}
L_1\sim\operatorname{Bernoulli}(\rho),\qquad
L_2\mid L_1\sim\operatorname{Bernoulli}\{s(\alpha+\beta L_1)\},
\quad s(t)=\sigmoid(t).
\label{eq:binarymodel}
\end{equation}
This directly models \(p(L_2\mid L_1)\). Under
\(q(L_1)q(L_2)\), let \(r_1=q(L_1=1)\) and \(r_2=q(L_2=1)\).
Since \(v^2=v\) for binary \(v\), the likelihood log-odds contribution
is \(b_k-a_k/2\). The child update is
\begin{equation}
\logit(r_2^{\rm new})=b_2-a_2/2+\alpha+\beta r_1.
\label{eq:binarychild}
\end{equation}
For the parent, write \(s_0=s(\alpha)\) and \(s_1=s(\alpha+\beta)\).
Its update contains feedback from the inferred child activation:
\begin{align}
\logit(r_1^{\rm new})
={}&b_1-a_1/2+\logit(\rho)\notag\\
&+r_2\log(s_1/s_0)+(1-r_2)\log\{(1-s_1)/(1-s_0)\}.
\label{eq:binaryparent}
\end{align}
If \(\beta>0\), a confidently active child tends to favor parent
activation. If \(\beta=0\), the child correction vanishes and both equations
reduce to ordinary independent Bernoulli-prior normal-means updates.

This example distinguishes three objects: the learned conditional prior
in \eqref{eq:binarymodel}, the posterior approximation \(r_1,r_2\), and the
posterior feedback required by that prior. A factorized q does not remove
the conditional relationship being learned in p.

\section{Keeping the feature update as ordinary ash}
Under the scalar-factorized approximation \eqref{eq:q}, set
\begin{align}
a^F_{jk}&=\sum_iw_{ij}\E_q[L_{ik}^2],\notag\\
b^F_{jk}&=\sum_iw_{ij}
\left(Y_{ij}-\sum_{h\ne k}\bar L_{ih}\bar F_{jh}\right)\bar L_{ik}.
\label{eq:Fstats}
\end{align}
Fit the existing ash normal-means model at
\(b^F_{jk}/a^F_{jk}\), with SE \((a^F_{jk})^{-1/2}\). Update its mixture
weights and posterior moments at each fitting iteration. Its configured
penalty and other options are preserved.

This is the original feature-side empirical Bayes update. Merely drawing F
from a prior fitted during initialization would not implement this learning
algorithm. No HMM or feature-side neural model is needed for this derivation.

\clearpage
\section{Noise, entropy and regularization}
\subsection{Observation noise}
Under \eqref{eq:q}, the expected squared residual is
\begin{align}
\E_q[(Y_{ij}-L_i^TF_j)^2]
={}&(Y_{ij}-\bar L_i^T\bar F_j)^2\notag\\
&+\sum_k\{\E[L_{ik}^2]\E[F_{jk}^2]-\bar L_{ik}^2\bar F_{jk}^2\}.
\label{eq:residual}
\end{align}
For unknown constant precision,
\begin{equation}
\tau^{\rm new}=\frac{|\Omega|}
{\sum_{ij\in\Omega}\E_q[(Y_{ij}-L_i^TF_j)^2]}.
\label{eq:noise}
\end{equation}
The corresponding row- or column-specific versions use the same sums
restricted to their groups. Known noise remains fixed.

\subsection{Entropy of the actual variational approximation}
For a continuous loading component in \eqref{eq:moments},
\begin{equation}
H_h=\tfrac12\log(2\pi V_h)
+\frac{\E_h[(v-m_h)^2]}{2V_h}-\E_h[D(v)]+\log I_{0h}.
\label{eq:entropy}
\end{equation}
The value-and-label entropy is
\(-\sum_hr_h\log r_h+\sum_{h>0}r_hH_h\). The atom has zero
within-component entropy. For F, use the ordinary EBNM identity to compute
its prior KL contribution. Thus the ELBO can be evaluated for the actual
q used by the algorithm, rather than a separate distribution fitted
afterward to sampler marginals.

When later coordinates or priors change, a previously updated q must stay
fixed until its next explicit update. A stored tilted density cannot silently
reevaluate itself using newly changed parent distributions or neural
parameters. Its representation must preserve the state that defined it.

\subsection{A consistent sparsity penalty}
First validate the unpenalized case. The existing neural spike regularizer
can then be included as
\begin{equation}
R(q,\theta)=\sum_{i,k}(\lambda_k-1)
\E_q\log\pi_{k0}(P_{ik}),\qquad \lambda_k\geq1.
\label{eq:penalty}
\end{equation}
The own penalty is constant in the candidate \(L_{ik}\), so it can be added outside
the log normalizer during \(\theta_k\) optimization. A child's penalty depends
on \(L_{ik}\) and must be included in \(D(v)\). Without latent parents these become
the fixed-covariate parameter penalties, preserving the ordinary penalized
cEBNM limit.

This is a regularized ELBO, or a bound for an augmented observation model;
it is not the unpenalized model's evidence. The ash global mixture-weight
penalty has its existing, separate semantics and should be preserved.

\clearpage
\section{The learning algorithm}
\note{\textbf{One cEBMF iteration.}
Retain the original residual-based \(L_k/F_k\) schedule. Extend only the loading
fit to account for uncertain parents and child feedback.}
\begin{enumerate}
\item Initialize \(L,F\) and their moments using the existing cEBMF scheme.
Preserve supplied factors and their scaling. There is no universal
statistical justification for thresholding initial loadings at 0.5.
\item For \(k=1,\ldots,K\):
  \begin{enumerate}[label=(\alph*)]
  \item Calculate \(a^L,b^L\) from current variational moments using
  \eqref{eq:stats}.
  \item If \(k\) has no latent parents or children, call its original cEBNM
  solver. Otherwise compute \eqref{eq:own}--\eqref{eq:child} and jointly
  fit the conditional prior and posterior through \eqref{eq:learning}.
  \item Store loading component probabilities, moments and the fixed
  variational representation needed for later expectations and entropy.
  \item Update \(F_k\) with the original ash fit using \eqref{eq:Fstats}, and
  refresh residuals.
  \end{enumerate}
\item Update unknown noise using the expected residuals and evaluate the
same variational objective.
\item Repeat until the chosen iteration limit or convergence criterion.
\end{enumerate}
The target remains \(\prod_kp_\theta(L_k\mid L_{<k})\). This variational
algorithm does not require a final full-matrix MCMC burn-in and sampling stage.

\subsection{Where numerical approximation remains}
Each integral over the target loading in \eqref{eq:integrals} is
one-dimensional. Adaptive quadrature or Gauss--Hermite integration around
the Gaussian base components can evaluate it. Expectations over several
uncertain parents can require separate quadrature or Monte Carlo integration.
Keep parent integration nodes or draws fixed during a local parameter fit,
and check results by increasing integration accuracy.

Sharp child constraints can make these calculations demanding. A practical
implementation must assess runtime and approximation error; the derivation
does not guarantee a fast neural optimizer or a predictive improvement.
When there are no latent edges, however, use the analytic cEBNM route and
avoid introducing numerical approximation unnecessarily.

\subsection{Compatibility and convergence requirements}
The no-edge route should preserve the old prior options, initialization,
RNG behavior, \(L_k/F_k\) schedule, noise updates and pruning policy. With a
hierarchy, deleting a factor changes child input identities and needs an
explicit graph/network policy. Fixing rank is reasonable for an initial
controlled comparison only when the comparator uses the same rank.

The exact variational argument gives nondecrease of the ELBO for exact
block fits. Existing approximate scalar solvers and finite neural updates
are not all exact maximizers; behavioral compatibility must not be confused
with an unconditional monotonicity theorem.

\clearpage
\section{If within-row posterior dependence is also required}
\label{sec:structured}
The same conditional-prior model can also be fitted with posterior covariance
between loading coordinates. For this additional inference choice, consider
\[
q(L,F)=\prod_iq_i(L_i)\prod_{j,k}q^F_{jk}(F_{jk}).
\]
The optimal row coordinate is
\begin{equation}
q_i^*(L_i)\propto p_\theta(L_i\mid X_i)
\exp\{-\tfrac12L_i^TH_iL_i+h_i^TL_i\},
\label{eq:rowblock}
\end{equation}
where
\[
H_i=\sum_jw_{ij}\E[F_jF_j^T],\qquad
h_i=\sum_jw_{ij}Y_{ij}\E[F_j].
\]
This retains the same conditional-prior model, but the feature numerator
must now use cross moments:
\begin{equation}
b^F_{jk}=\sum_iw_{ij}\left[
Y_{ij}\E[L_{ik}]-\sum_{h\ne k}\E[L_{ik}L_{ih}]\E[F_{jh}]\right].
\label{eq:cross}
\end{equation}
Relative to \eqref{eq:Fstats}, the extra term is
\[
-\sum_iw_{ij}\sum_{h\ne k}\Cov(L_{ik},L_{ih})\bar F_{jh}.
\]
Thus one cannot sample a correlated L block and pass only marginal means
and variances into an unchanged mean-residual feature update.

Full row blocks do not automatically recover mean-field cEBMF when prior
edges vanish: the likelihood can still correlate loading coordinates.
For exact reduction, define blocks by connected components of the prior
graph and split them into singleton columns in the no-edge case.
Singleton blocks recover ordinary cEBNM; larger blocks require a
multivariate profile objective.

\section{Verification and status}
The repository script
\path{output/review_joint_l/verify_replacement_derivation.py} checks:
\begin{itemize}
\item Gaussian-mixture probabilities, moments and log normalizers in the
no-parent/no-child limit: maximum absolute error \(2.0\times10^{-15}\).
\item Independent and dependent binary coordinates against exhaustive
enumeration: maximum probability error \(1.2\times10^{-16}\), with
nonnegative exact coordinate ELBO gains.
\item The profile objective against the directly optimized enumerated ELBO
at four parameter settings: their difference is constant to
\(6.7\times10^{-16}\).
\item The structured feature numerator: a correlated binary example gives
the correct value 0; multiplying marginal means incorrectly gives 0.25.
\end{itemize}
The implementation in \path{cebmf/_conditional.py} additionally has analytic
profile-gradient, no-edge compatibility and missing-modality regression tests.
Its integration approximations and finite training still require sensitivity
checks and comparative empirical evaluation.

\clearpage
\section{Adversarial numerical review and device handling}
\label{sec:numerical-audit}
\textit{Review added 16 September 2026.} The target remains
\(p(L_k\mid L_{<k})\). Removing latent-covariate edges gives the same
normal-means empirical Bayes subproblem; the implementation dispatches this
case to the ordinary cEBMF solver. This algebraic statement does not assert
that every finite neural optimization attains the subproblem optimum.

\subsection{A correct formula can be an unstable implementation}
Completing a Gaussian square using its natural parameters subtracts terms
of order \(\mu^2/s^2\). In float32 this failed for a prior with spike weight
0.2 and slab \(\N(10,10^{-8})\), with likelihood statistics \(a=2,b=3\).
The old numerical profile returned a posterior mean near 0.0125; the
independent analytic posterior is approximately \(1.59\times10^{-29}\).
This is cancellation error, not Monte Carlo variability.

For each continuous component, compute instead
\begin{align}
t&=\E[s(P)^{-2}], &
c&=\frac{\E[\mu(P)s(P)^{-2}]}{t},\\
h&=\E\log\pi(P)-\frac12\E\!\left[
 \log(2\pi s(P)^2)+\frac{(\mu(P)-c)^2}{s(P)^2}\right].
\end{align}
Then the expected log prior is exactly
\[
A(v)=h-\tfrac12t(v-c)^2.
\]
The spread in this expression is evaluated as an average of squared
centered differences, never as a difference of large second moments.
The height \(h\) retains the parameter-dependent normalization of
\(\exp(\E\log g)\); dropping it would change empirical Bayes learning.

Write \(\alpha=a+t\), \(r=t/\alpha\). The likelihood-completed Gaussian
has mean \(m=rc+b/\alpha\), variance \(V=1/\alpha\), and log mass
\begin{equation}
B=h+r(-ac^2/2+bc)+\frac{b^2}{2\alpha}
  +\frac12\log\frac{2\pi}{\alpha}.
\end{equation}
Its tilted contribution is \(\exp(B)\E_{N(m,V)}\exp D(v)\).
For the zero atom the corresponding log mass before the child tilt is
\(\E\log\pi_0(P)\). These are algebraic rearrangements of the earlier
derivation, not a different statistical model.

At a standard-normal quadrature node \(x_j\), evaluate the slab log density
as \(-\tfrac12\{\log(2\pi V)+x_j^2\}\). Combining it with \(B+D(v)-\log Z\)
gives the continuous-and-label log posterior used for entropy. This avoids
another unstable natural-parameter polynomial. For child feedback use
\[
\E[(L_c-\mu(P))^2\mid Z_c=h]
=\operatorname{Var}(L_c\mid Z_c=h)
 +(\E[L_c\mid Z_c=h]-\mu(P))^2,
\]
computing the within-component variance directly from centered nodes.

\subsection{What was checked, and what was not established}
Independent checks include exhaustive binary enumeration, analytic mixture
posteriors (including zero likelihood precision), and a dense trapezoidal
integral for a nonlinear logistic child tilt. Evidence, moments, entropy,
and parameter gradients agree at the stated test tolerances. Direct child
integration also checks the sparsity penalty and missing-branch mask.
The ordinary spike--normal posterior now uses log-domain component odds;
normal-mixture likelihoods are no longer silently clipped in the tails.

These tests do not prove convergence of finite quadrature for every neural
prior. Multiple-parent expectations use a finite fixed Sobol approximation;
the stored quadrature representation also approximates each continuous
coordinate. Local profile checkpoint selection is not a proof of global
ELBO ascent under those approximations. The posterior is still factorized
across loading coordinates. Better prior flexibility can lose predictive
accuracy through optimization, approximation, or overfitting.

The saved tree benchmark predates these numerical fixes. Inspecting its
checkpoints confirms float32 observations and loadings; the earlier claim
of float64 fitting was incorrect. The narrow-slab bug could affect those
fits, but attributing the performance gap to it requires matched reruns.
Original source hashes and measurements are retained and have not been
relabeled as results of the revised implementation.

\subsection{CUDA contract and its validation limit}
Loading moments, networks, quadrature states, optimizer state, residuals,
and objective histories are kept on the selected device. Conditional
checkpoint selection and mixture-EM convergence use device tensors.
Observation indices and prediction tables are cached. Sobol tables are
constructed once on CPU and transferred at graph setup. Ordinary CUDA
neural solvers batch with device indices.

ASH's automatically sized scale grid needs one host scalar for its length;
its amplitudes stay on device. Supplying a fixed \texttt{scales} tensor to
the ash prior removes that dynamic-size decision and still refits weights
on every sweep. Fixed-grid choice is a modeling choice, not an automatically
equivalent replacement for an adaptive grid. Rank pruning, boundary input
validation, plotting, and explicit reporting may also synchronize.

The development environment has PyTorch 2.11.0+cpu and no CUDA runtime.
CUDA-only tests cover complete warmed sweeps, partially paired modalities,
optimizer-state residency, CPU/GPU numerical agreement, profiler transfers,
and retained memory. They are skipped here, not reported as GPU passes.
Hardware validation remains required before making a measured GPU speed or
memory claim. The deprecated experimental sampler and optional L-BFGS route
are outside the synchronization-free conditional/ash-EM contract.


\vspace{3pt}
\textbf{Reference.} Denault et al. (2025),
\emph{Covariate-moderated Empirical Bayes Matrix Factorization},
Section 3.3 and its variational derivation.
\href{https://arxiv.org/abs/2505.11639}{arXiv:2505.11639}.
The conditional-prior extension and reduction above are derived for this project.
\end{document}
